%% 10-conflict-resolution-and-governance.tex -- generated by tools/build_from_docx.py; edit freely

\hypertarget{conflict-resolution-and-governance}{%
\section{Conflict Resolution and Governance}\label{conflict-resolution-and-governance}}

Conflict resolution is a contribution of Work Package 4, and this
chapter does not design conflict-resolution algorithms. What it
specifies is the contract by which an agent interacts with the conflict
fabric, and the guardrails that make enforcement safe and deterministic.
The division is deliberate: Work Package 4 decides, and Task 2.3 binds.
The architecture work names the open gap directly, observing that
agentic conflict handling lacks operational guardrails for safe,
deterministic enforcement in real networks, and that gap is what this
chapter addresses.

The distinction matters practically. Any resolution algorithm the
project adopts must be substitutable without a single admitted agent
changing, which is only possible if the surface agents see is fixed and
narrow.

\hypertarget{what-is-adopted-as-given}{%
\subsection{What is Adopted as Given}\label{what-is-adopted-as-given}}

The conflict taxonomy is taken as supplied. By scope, a Type 1 conflict
is intra-domain and resolved at domain level over network parameters and
key performance indicators; a Type 2 conflict is inter-domain, between a
domain and the inter-domain management function; a Type 3 conflict
arises from inter-domain management decisions such as activating
additional services, and is resolved by negotiation between the domain
and inter-domain management functions.

By detectability, a direct conflict involves a shared control parameter.
An indirect conflict involves distinct control parameters affecting a
shared key performance indicator, found by clustering control parameters
over an association matrix. An implicit conflict is detected reactively,
through unexpected degradation of a key performance indicator, which
then updates the control-parameter clusters.

The resolution machinery is likewise adopted rather than redesigned: the
conflict detector with its extraction, checking, clustering and
forwarding modules; the conflict resolver applying priority-based,
service-level-agreement-based or fairness-based policies; network
digital twin pre-simulation of candidate actions before deployment; the
distributed variant emitting conflict events from a smart contract; and
the generator producing inter-domain contracts whose logic is not known
in advance. Transport is the management bus, with a status topic
carrying detected conflicts and their types, and request and response
topics carrying negotiation between the domain and inter-domain
management functions.

\hypertarget{the-agent-facing-contract}{%
\subsection{The Agent-Facing Contract}\label{the-agent-facing-contract}}

Task 2.3 defines exactly four interactions. This is the whole surface an
agent sees, and any resolution algorithm may sit behind it.

\hypertarget{register}{%
\subsubsection{Register}\label{register}}

Each agent declares its control parameters and its target key
performance indicators. These already exist in the descriptor as the
control scope and action space, which are the control parameters, and
the declared objectives, which are the key performance indicators.

Registration is performed by the platform from the admitted descriptor
and is never self-reported at runtime. This is what keeps the
control-parameter and key-performance-indicator matrices from
disagreeing with the authority the agent was actually granted. Detection
of direct and indirect conflicts is a matrix operation over that
membership; were registration voluntary or mutable at runtime, an agent
could acquire authority invisible to the detector, which defeats
detection rather than resolving conflict.

\hypertarget{submit}{%
\subsubsection{Submit}\label{submit}}

At intent egress the Decision Engine's runtime submits the intent for
mediation. The submission is the intent envelope of Chapter 6, from
which the fabric reads the identity of the submitting instance and its
generation, the managed objects and control parameters the intent would
change, the objective it pursues, the outcome it seeks and its validity
window, together with an urgency claim and a candidate set where
present. Table 10-1 states what the fabric reads and how it uses each
element.

\begin{wptable}{10-1}

\begin{longtable}[]{@{}
  >{\raggedright\arraybackslash}p{(\columnwidth - 4\tabcolsep) * \real{0.2501}}
  >{\raggedright\arraybackslash}p{(\columnwidth - 4\tabcolsep) * \real{0.0750}}
  >{\raggedright\arraybackslash}p{(\columnwidth - 4\tabcolsep) * \real{0.6750}}@{}}
\caption{Content of a submission to the conflict fabric}\tabularnewline
\toprule\noalign{}
\begin{minipage}[b]{\linewidth}\raggedright
\textbf{Field of the intent envelope}
\end{minipage} & \begin{minipage}[b]{\linewidth}\raggedright
\textbf{Req.}
\end{minipage} & \begin{minipage}[b]{\linewidth}\raggedright
\textbf{Use by the fabric}
\end{minipage} \\
\midrule\noalign{}
\endfirsthead
\toprule\noalign{}
\begin{minipage}[b]{\linewidth}\raggedright
\textbf{Field of the intent envelope}
\end{minipage} & \begin{minipage}[b]{\linewidth}\raggedright
\textbf{Req.}
\end{minipage} & \begin{minipage}[b]{\linewidth}\raggedright
\textbf{Use by the fabric}
\end{minipage} \\
\midrule\noalign{}
\endhead
\bottomrule\noalign{}
\endlastfoot
\textbf{intent\_\allowbreak{}id} & M & Correlates the submission with
its verdict, its token and any execution envelope. \\
\textbf{agent\_\allowbreak{}instance\_\allowbreak{}id} & M & Identity of
the submitting instance, with the generation the runtime attaches;
resolved against the registry, never trusted from the message. \\
\textbf{scope, desired\_\allowbreak{}outcome} & M & Managed objects,
action and parameter values sought. The control parameters affected are
derived from them, and a submission exceeding the registered control
scope is a conformance fault. \\
\textbf{objective\_\allowbreak{}ref} & M & Objective pursued; its KPI is
the one registered for the agent. It is what lets the fabric arbitrate
between competing objectives as well as over shared parameters. \\
\textbf{validity\_\allowbreak{}window} & M & Period beyond which the
intent is stale and must not be actuated. \\
\textbf{urgency} & O & Orders the agent's own intents within its
registered priority band; never raises the band. \\
\textbf{candidates{[}{]}} & O & Alternative outcomes meeting the same
objective; a modify verdict may name the one the fabric permits. \\
\textbf{rationale} & M & Carried for audit, and for arbitration where
the resolution policy uses it; never a source of authority. \\
\end{longtable}

\end{wptable}

Arbitration is therefore not confined to contention for a resource. Two
intents on distinct control parameters conflict when the key performance
indicators they move are objectives pulling against each other, as when
one agent maximizes an indicator that another minimizes or trades
against another's target, and the fabric resolves such an objective
conflict under the same policies and with the same four verdicts as a
contest for a shared parameter. The objective reference in the
submission is what makes this possible: without it the fabric could see
which parameters an intent would change, but not which goal the change
serves.

\hypertarget{verdict}{%
\subsubsection{Verdict}\label{verdict}}

The fabric returns exactly one of four verdicts, with a reason code and
the identity of the resolving authority. Allow permits the intent as
submitted. Modify carries a bounded alternative, re-validated against
the agent's own action space and treated as deny where it falls outside.
Defer carries a retry-after, and the agent must not actuate in the
interim. Deny is terminal for that intent and re-enters the Decision
Engine as evidence. Table 10-2 states the consequence of each verdict
for the submitting agent.

\begin{wptable}{10-2}

\begin{longtable}[]{@{}
  >{\raggedright\arraybackslash}p{(\columnwidth - 4\tabcolsep) * \real{0.1375}}
  >{\raggedright\arraybackslash}p{(\columnwidth - 4\tabcolsep) * \real{0.3375}}
  >{\raggedright\arraybackslash}p{(\columnwidth - 4\tabcolsep) * \real{0.5250}}@{}}
\caption{The four verdicts and their consequences for the
agent}\tabularnewline
\toprule\noalign{}
\begin{minipage}[b]{\linewidth}\raggedright
\textbf{Verdict}
\end{minipage} & \begin{minipage}[b]{\linewidth}\raggedright
\textbf{Meaning}
\end{minipage} & \begin{minipage}[b]{\linewidth}\raggedright
\textbf{Consequence for the agent}
\end{minipage} \\
\midrule\noalign{}
\endfirsthead
\toprule\noalign{}
\begin{minipage}[b]{\linewidth}\raggedright
\textbf{Verdict}
\end{minipage} & \begin{minipage}[b]{\linewidth}\raggedright
\textbf{Meaning}
\end{minipage} & \begin{minipage}[b]{\linewidth}\raggedright
\textbf{Consequence for the agent}
\end{minipage} \\
\midrule\noalign{}
\endhead
\bottomrule\noalign{}
\endlastfoot
\textbf{allow} & The intent proceeds as submitted. & Token issued;
actuation permitted within the validity window. \\
\textbf{modify} & A bounded alternative is returned. & Re-validated
against the agent\textquotesingle s own action space; out of bounds is
treated as deny. \\
\textbf{defer} & No decision yet; a retry-after is returned. & The agent
must not actuate in the interim. \\
\textbf{deny} & The intent is refused. & Terminal for that intent;
re-enters the DE as evidence. \\
\end{longtable}

\end{wptable}

Where a Decision Engine can meet its objective by more than one action,
the submission may carry a candidate set rather than a single proposed
intent, and the fabric may return modify naming the candidate it
permits. This is negotiation-shaped in one round trip and requires no
interaction beyond the four, which is why it is the form the contract
adopts. Should Work Package 4 later require a genuinely multi-round
exchange, the available extension is a bounded one, in which the fabric
defers with a counter-proposal and the agent resubmits against it under
a hard round cap; that introduces no new interaction either, but it
makes mediation latency variable and would require the effective latency
class to be computed against the cap rather than a single hop. An
agent-initiated negotiation is deliberately not provided: an agent that
bargains with the fabric rather than requesting permission from it is
the point at which self-approval re-enters by another route.

This vocabulary is a normalization rather than a new algorithm.
Priority, arbitration, negotiation and hierarchical policies all reduce
to these four verdicts. Agents are built by different partners against
different algorithms, and a fixed minimal verdict vocabulary is
precisely what allows the resolution algorithm to be replaced without
touching a single agent.

\hypertarget{notify}{%
\subsubsection{Notify}\label{notify}}

The fabric may asynchronously notify an agent of an implicit conflict
attributed to it. A notification is not a verdict on a pending intent;
it re-enters the agent as evidence and may trigger an evaluation-led
downgrade. Notification also carries the screening-basis changes of
Section 10.3.1, by which an agent learns that a parameter it held
exclusively has become shared, or that an indicator its parameters move
has become another agent's objective. Implicit conflicts are detected
only after damage appears in the key performance indicators, so they
cannot be handled synchronously, but they must not be silently discarded
either.

\hypertarget{mediation-is-synchronous-and-fail-closed}{%
\subsection{Mediation is Synchronous and Fail-Closed}\label{mediation-is-synchronous-and-fail-closed}}

An intent without a verdict is not actuatable. Where mediation times
out, the declared unresolved-conflict behavior applies, and never an
implicit allow. The mediation binding is verified at the activation gate
and the intent-egress gate consults it every cycle.

Asynchronous mediation is self-approval with a delay. If an agent could
actuate while a verdict was outstanding, a conflict over a shared
control parameter would be resolved only after both actions had already
landed on it, which is detection rather than resolution.

The cost is acknowledged: mediation sits in the control path and
consumes loop budget. This is why the effective latency class is
computed with the mediation hop included. An agent claiming a loop
budget too tight to mediate fails admission rather than silently
bypassing the fabric.

\hypertarget{screening-before-submission}{%
\subsubsection{Screening before submission}\label{screening-before-submission}}

Mediation is not invoked on every decision. The Decision Engine screens
each candidate against cached bounds, resource-sharing metadata and
control-parameter coupling, and posts a resolution request only where
the action touches a shared resource, exceeds a bound, changes a
parameter controlled by another agent, moves a key performance indicator
that another agent has registered as its objective, or creates an
unresolved cross-domain effect.

Screening is a filter on which candidates require a verdict, not a
self-granted approval. An action requiring authority still cannot
proceed without one, and the cached bounds against which screening tests
are themselves previously granted authority. This is the mechanism by
which mediation remains compatible with tight loop budgets, and it is
what keeps the exclusivity case of Chapter 6 from becoming an exception
the agent itself decides.

What makes screening safe over time is that the basis it tests against
is maintained by the fabric rather than held as a snapshot by the agent.
Because registration records which agents hold which control parameters
and which key performance indicators each pursues, the fabric knows the
moment a parameter ceases to be exclusive, and the moment an agent's
parameters come to move another agent's objective. When a second agent
registers against a parameter one agent previously held alone, or
registers an objective on an indicator that the incumbent's parameters
affect by the association matrix, the fabric notifies the incumbent, and
the notification carries one of two things: either that the parameter is
now contested and actions touching it must be submitted from this point,
or a revised set of bounds within which each agent may continue to
operate without submission. The screening basis is invalidated by the
fabric at the moment it changes, and not left to expire.

This matters because expiry alone would leave a window in which an agent
screens correctly against a world that no longer exists. Driving
invalidation from registration closes that window, and it places the
obligation where the knowledge is: only the fabric can see a second
registration arriving, and the incumbent agent has no way to learn of it
by itself.

\hypertarget{bounded-approval-and-stored-validity}{%
\subsubsection{Bounded approval and stored validity}\label{bounded-approval-and-stored-validity}}

A verdict may return a bounded approval together with its validity
period, a mitigation and an audit reference. The Decision Engine stores
the returned bounds for their validity period and continues autonomous
operation within them, without a submission per cycle.

This is the action envelope of Section 3.9.2, and it is what makes
mediation compatible with the tightest latency classes: the round trip
is paid once for a window of operation rather than once per decision.

\hypertarget{the-conflict-resolution-token}{%
\subsubsection{The conflict-resolution token}\label{the-conflict-resolution-token}}

The verdict carries a token, and the Actuator validates that token at
the last mile alongside semantic authority, target scope and action
bounds. An intent arriving at actuation without a valid, in-validity
token is rejected there.

The token is what turns ``authority was granted upstream'' from an
assumption into an artifact checked at the point of execution. It is the
concrete form of the last-mile re-validation specified in Chapter 7.

\hypertarget{placement-off-the-fast-path}{%
\subsubsection{Placement off the fast path}\label{placement-off-the-fast-path}}

Conflict resolution and evaluation use explicit exceptional bindings and
do not sit on the ordinary fast path. Screening, cached bounds and
stored validity are what keep the ordinary cycle clear of the fabric;
the fabric is reached on the exceptional path, when screening says a
candidate needs it.

\hypertarget{escalation}{%
\subsection{Escalation}\label{escalation}}

Type 1 conflicts resolve locally at the domain management function. Type
2 escalate from the domain to the inter-domain management function. Type
3 are negotiated between the two over the request and response topics.

An agent never initiates escalation and never addresses the inter-domain
management function. It submits locally and receives one verdict.
Escalation requires a view across agents and domains that no single
agent possesses; letting agents escalate would recreate the global-brain
topology the framework rejects, and would let a partner's agent route
around the authority of its own domain.

This has a consequence for federated agents. The per-domain members of
one logical agent are mediated independently, each in its own domain.
Cross-domain incoherence between them surfaces as a Type 2 conflict,
which means the fabric is precisely the mechanism that makes federation
safe without cross-domain subscription between the members.

\hypertarget{governance}{%
\subsection{Governance}\label{governance}}

\hypertarget{the-registry-is-the-single-source-of-authority-truth}{%
\subsubsection{The registry is the single source of authority truth}\label{the-registry-is-the-single-source-of-authority-truth}}

The domain management function holds, per agent instance: identity,
generation, authority state, control scope, action space, registered
control-parameter and key-performance-indicator sets, declared failure
policy, conflict policy and priority rules. The fabric resolves against
the registry and never against claims carried in a submitted message.
Table 10-3 states what the registry holds per instance.

\begin{wptable}{10-3}

\begin{longtable}[]{@{}
  >{\raggedright\arraybackslash}p{(\columnwidth - 2\tabcolsep) * \real{0.3499}}
  >{\raggedright\arraybackslash}p{(\columnwidth - 2\tabcolsep) * \real{0.6501}}@{}}
\caption{Per-instance registry content held by the DMO}\tabularnewline
\toprule\noalign{}
\begin{minipage}[b]{\linewidth}\raggedright
\textbf{Registry content}
\end{minipage} & \begin{minipage}[b]{\linewidth}\raggedright
\textbf{Role in mediation}
\end{minipage} \\
\midrule\noalign{}
\endfirsthead
\toprule\noalign{}
\begin{minipage}[b]{\linewidth}\raggedright
\textbf{Registry content}
\end{minipage} & \begin{minipage}[b]{\linewidth}\raggedright
\textbf{Role in mediation}
\end{minipage} \\
\midrule\noalign{}
\endhead
\bottomrule\noalign{}
\endlastfoot
\textbf{Identity and generation} & Fences retired instances at the
fabric. \\
\textbf{Authority state} & Consulted per cycle at intent egress. \\
\textbf{Control scope and action space} & Bounds every submission;
excess is a conformance fault. \\
\textbf{Registered CP and KPI sets} & Membership over which direct and
indirect detection operate. \\
\textbf{Declared failure policy} & Applied on timeout and on unresolved
conflict. \\
\textbf{Conflict policy and priority rules} & Effective arbitration
priority, set at admission from the service tier. \\
\end{longtable}

\end{wptable}

A submission whose control parameters exceed the registered control
scope is rejected as a conformance fault rather than arbitrated. Out of
scope is not a conflict; it is a violation. Arbitration presumes that
contending parties are each entitled to what they ask for, and an agent
exceeding its scope is not a party to a conflict. Treating violation as
conflict would let an agent enlarge its authority simply by asking for
more.

\hypertarget{priority-claims-are-bounded-by-the-registry}{%
\subsubsection{Priority claims are bounded by the registry}\label{priority-claims-are-bounded-by-the-registry}}

An agent may express urgency, but its effective arbitration priority is
the one registered at admission from the service tier. Self-declared
urgency may order intents only within the agent's own registered
priority band.

Without this bound, every partner maximizes its own urgency field and
priority-based resolution collapses into first-come-first-served. The
bound is what keeps a priority policy meaningful under multi-tenant
contribution, which is the exact condition the platform has to survive.

\hypertarget{digital-twin-validation-is-resolver-side}{%
\subsubsection{Digital twin validation is resolver-side}\label{digital-twin-validation-is-resolver-side}}

Where the conflict management function uses a network digital twin to
simulate candidate actions before deployment, that simulation belongs to
the resolver. An agent may not assert that its action was
twin-validated, and no authority is granted on the strength of such a
claim.

If agents could self-certify twin validation, the assurance would be
exactly as trustworthy as the least rigorous partner implementation.
Held resolver-side it is uniform across every tenant. Twin validation
strengthens a verdict; it never strengthens an agent's authority state.

\hypertarget{conflict-outcomes-as-evidence}{%
\subsection{Conflict Outcomes as Evidence}\label{conflict-outcomes-as-evidence}}

Per agent and per service, the platform collects the verdict
distribution across allow, modify, defer and deny; the unresolved
conflict count; the escalation count; implicit-conflict attributions;
and oscillation indicators.

These feed three consumers. They enter the requester-facing assurance
report of Chapter 9. They feed evaluation-led authority downgrade. And
they are the measurement point for the project's conflict-penalty
reduction target, since the same counters quantify both agent behavior
and fabric effectiveness.

A high deny rate or a high implicit-attribution rate is the earliest
evidence that an agent is mis-specified or unsafe, and it is available
long before that agent's actions damage a service key performance
indicator. An agent repeatedly denied is telling the platform something
about itself that no self-report would disclose. The counters are
platform-emitted and never agent-reported, for the same reason as
everywhere else in this framework. Table 10-4 states the indicators and
the consumers of each.

\begin{wptable}{10-4}

\begin{longtable}[]{@{}
  >{\raggedright\arraybackslash}p{(\columnwidth - 4\tabcolsep) * \real{0.2250}}
  >{\raggedright\arraybackslash}p{(\columnwidth - 4\tabcolsep) * \real{0.5000}}
  >{\raggedright\arraybackslash}p{(\columnwidth - 4\tabcolsep) * \real{0.2750}}@{}}
\caption{Conflict indicators and their consumers}\tabularnewline
\toprule\noalign{}
\begin{minipage}[b]{\linewidth}\raggedright
\textbf{Indicator}
\end{minipage} & \begin{minipage}[b]{\linewidth}\raggedright
\textbf{Definition}
\end{minipage} & \begin{minipage}[b]{\linewidth}\raggedright
\textbf{Consumers}
\end{minipage} \\
\midrule\noalign{}
\endfirsthead
\toprule\noalign{}
\begin{minipage}[b]{\linewidth}\raggedright
\textbf{Indicator}
\end{minipage} & \begin{minipage}[b]{\linewidth}\raggedright
\textbf{Definition}
\end{minipage} & \begin{minipage}[b]{\linewidth}\raggedright
\textbf{Consumers}
\end{minipage} \\
\midrule\noalign{}
\endhead
\bottomrule\noalign{}
\endlastfoot
\textbf{Verdict distribution} & allow / modify / defer / deny counts per
agent and per service. & Assurance; evaluation; KPI\#21 \\
\textbf{Unresolved conflict count} & Submissions terminating without a
verdict. & Assurance; evaluation \\
\textbf{Escalation count} & Conflicts escalated from DMO to IDMO. &
Assurance; KPI\#21 \\
\textbf{Implicit-conflict attributions} & Conflicts attributed to the
agent through the KPI-degradation path. & Evaluation; authority
downgrade \\
\textbf{Oscillation indicators} & Repeated reversal over the same
control parameter. & Evaluation; authority downgrade \\
\end{longtable}

\end{wptable}

\hypertarget{requirements-7}{%
\subsection{Requirements}\label{requirements-7}}

\begin{requirement}{Req-T23-CG-01}{MUST}{Registration is performed by the platform}
\reqfield{Description / rationale}{%
The platform shall register each agent\textquotesingle s control
parameters and target KPIs from the admitted descriptor. Registration
shall not be self-reported at runtime and shall not be mutable by the
agent, so that the CP and KPI matrices cannot disagree with the
authority actually granted.}
\reqfield{Dependency}{%
Req-T23-LC-01}
\reqfield{Traceability (Forward)}{%
\begin{itemize}
\tightlist
\item
  Derived from the descriptor at admission.
\end{itemize}}
\reqfield{UNITY-6G components}{%
\begin{itemize}
\tightlist
\item
  DMO
\item
  Admission (IDMO/DMO)
\end{itemize}}
\reqfield{Linked KPIs}{%
\begin{itemize}
\tightlist
\item
  KPI\#21: Reduce the performance penalty due to conflicts by 20\% by
  means of DLT, hierarchical algorithms and policy-based conflict
  resolution schemes.
\item
  KPI\#36: Guarantee 99.99999\% service availability and reliability
  based on effective monitoring and LCM.
\end{itemize}}
\end{requirement}

\begin{requirement}{Req-T23-CG-02}{MUST}{The agent-facing surface is exactly four interactions}
\reqfield{Description / rationale}{%
The conflict fabric shall expose to agents exactly four interactions:
register, submit, verdict and notify. No resolution algorithm shall
require an agent-visible interaction outside this surface. A submission
may carry a candidate set, and a modify verdict may name the permitted
candidate; the fabric shall provide no agent-initiated negotiation.}
\reqfield{Dependency}{%
Req-T23-DE-03}
\reqfield{Traceability (Forward)}{%
\begin{itemize}
\tightlist
\item
  Surface fixed at the conflict fabric binding.
\end{itemize}}
\reqfield{UNITY-6G components}{%
\begin{itemize}
\tightlist
\item
  Conflict fabric
\item
  DMO
\item
  WP4 resolver
\end{itemize}}
\reqfield{Linked KPIs}{%
\begin{itemize}
\tightlist
\item
  KPI\#21: Reduce the performance penalty due to conflicts by 20\% by
  means of DLT, hierarchical algorithms and policy-based conflict
  resolution schemes.
\item
  KPI\#1: Deliver four distributed management entities (MS, AE, DE and
  ACT) deployable in each heterogeneous domain.
\end{itemize}}
\end{requirement}

\begin{requirement}{Req-T23-CG-03}{MUST}{The verdict vocabulary is fixed}
\reqfield{Description / rationale}{%
The fabric shall return exactly one of allow, modify, defer or deny,
with a reason code and the identity of the resolving authority. modify
shall carry a bounded alternative; defer shall carry a retry-after; deny
shall be terminal for that intent and shall re-enter the DE as evidence.}
\reqfield{Dependency}{%
Req-T23-CG-02}
\reqfield{Traceability (Forward)}{%
\begin{itemize}
\tightlist
\item
  Verdict schema enforced at the fabric binding.
\end{itemize}}
\reqfield{UNITY-6G components}{%
\begin{itemize}
\tightlist
\item
  Conflict fabric
\item
  DE runtime
\end{itemize}}
\reqfield{Linked KPIs}{%
\begin{itemize}
\tightlist
\item
  KPI\#21: Reduce the performance penalty due to conflicts by 20\% by
  means of DLT, hierarchical algorithms and policy-based conflict
  resolution schemes.
\item
  KPI\#36: Guarantee 99.99999\% service availability and reliability
  based on effective monitoring and LCM.
\end{itemize}}
\end{requirement}

\begin{requirement}{Req-T23-CG-04}{MUST}{Mediation is synchronous and fail-closed}
\reqfield{Description / rationale}{%
An intent without a verdict shall not be actuatable. On mediation
timeout the declared conflict\_\allowbreak{}unresolved behavior shall
apply, and shall never be an implicit allow. The mediation binding shall
be verified at the activation gate and consulted at intent egress every
cycle.}
\reqfield{Dependency}{%
Req-T23-DE-02}
\reqfield{Traceability (Forward)}{%
\begin{itemize}
\tightlist
\item
  Verified at activation gate check 7.
\item
  Consulted per cycle at intent egress.
\end{itemize}}
\reqfield{UNITY-6G components}{%
\begin{itemize}
\tightlist
\item
  Conflict fabric
\item
  DMO
\item
  Platform runtime
\end{itemize}}
\reqfield{Linked KPIs}{%
\begin{itemize}
\tightlist
\item
  KPI\#21: Reduce the performance penalty due to conflicts by 20\% by
  means of DLT, hierarchical algorithms and policy-based conflict
  resolution schemes.
\item
  KPI\#36: Guarantee 99.99999\% service availability and reliability
  based on effective monitoring and LCM.
\end{itemize}}
\end{requirement}

\begin{requirement}{Req-T23-CG-05}{MUST}{Screening selects candidates, it does not approve them}
\reqfield{Description / rationale}{%
The DE shall screen each candidate against cached bounds,
resource-sharing metadata and control-parameter coupling, and shall
submit where the action touches a shared resource, exceeds a bound,
changes a parameter controlled by another agent, moves a key performance
indicator registered as another agent\textquotesingle s objective, or
creates an unresolved cross-domain effect. Screening shall never
constitute approval of an action that requires authority.}
\reqfield{Dependency}{%
Req-T23-CG-04}
\reqfield{Traceability (Forward)}{%
\begin{itemize}
\tightlist
\item
  Screening criteria declared and checked at admission.
\end{itemize}}
\reqfield{UNITY-6G components}{%
\begin{itemize}
\tightlist
\item
  DE runtime
\item
  Conflict fabric
\end{itemize}}
\reqfield{Linked KPIs}{%
\begin{itemize}
\tightlist
\item
  KPI\#4: Real-time reconfiguration of IEEE 802.11 network domain in
  range smaller than 10 ms from taking the decision.
\item
  KPI\#21: Reduce the performance penalty due to conflicts by 20\% by
  means of DLT, hierarchical algorithms and policy-based conflict
  resolution schemes.
\item
  KPI\#23: Reduction of communications overhead by a factor of 10 while
  maintaining the SLA requirements.
\end{itemize}}
\end{requirement}

\begin{requirement}{Req-T23-CG-13}{MUST}{The fabric invalidates the screening basis on change}
\reqfield{Description / rationale}{%
Where registration causes a control parameter previously held by one
agent to become shared, registers an objective on a key performance
indicator that another agent\textquotesingle s control parameters
affect, or otherwise changes the coupling on which screening depends,
the fabric shall notify every affected agent. The notification shall
state either that actions touching the parameter must be submitted from
that point, or the revised bounds within which each affected agent may
continue without submission. The screening basis shall be invalidated by
the fabric at the moment it changes and shall not be left to expire.}
\reqfield{Dependency}{%
Req-T23-CG-05}
\reqfield{Traceability (Forward)}{%
\begin{itemize}
\tightlist
\item
  Invalidation driven by registration, not by cache expiry.
\item
  Notification delivered over the notify interaction.
\end{itemize}}
\reqfield{UNITY-6G components}{%
\begin{itemize}
\tightlist
\item
  Conflict fabric
\item
  DMO
\end{itemize}}
\reqfield{Linked KPIs}{%
\begin{itemize}
\tightlist
\item
  KPI\#21: Reduce the performance penalty due to conflicts by 20\% by
  means of DLT, hierarchical algorithms and policy-based conflict
  resolution schemes.
\item
  KPI\#36: Guarantee 99.99999\% service availability and reliability
  based on effective monitoring and LCM.
\end{itemize}}
\end{requirement}

\begin{requirement}{Req-T23-CG-06}{MAY}{Bounded approval carries stored validity}
\reqfield{Description / rationale}{%
A verdict may return a bounded approval with a validity period, a
mitigation and an audit reference. The DE may store the returned bounds
and continue autonomous operation within them for the validity period,
without submitting per cycle.}
\reqfield{Dependency}{%
Req-T23-CG-03}
\reqfield{Traceability (Forward)}{%
\begin{itemize}
\tightlist
\item
  Validity period enforced by the DE runtime and re-checked at ACT.
\end{itemize}}
\reqfield{UNITY-6G components}{%
\begin{itemize}
\tightlist
\item
  DE runtime
\item
  Conflict fabric
\end{itemize}}
\reqfield{Linked KPIs}{%
\begin{itemize}
\tightlist
\item
  KPI\#4: Real-time reconfiguration of IEEE 802.11 network domain in
  range smaller than 10 ms from taking the decision.
\item
  KPI\#23: Reduction of communications overhead by a factor of 10 while
  maintaining the SLA requirements.
\end{itemize}}
\end{requirement}

\begin{requirement}{Req-T23-CG-07}{MUST}{The conflict-resolution token is checked at actuation}
\reqfield{Description / rationale}{%
A verdict permitting action shall carry a conflict-resolution token. The
ACT shall validate the token alongside semantic authority, target scope
and action bounds at the last mile, and shall reject an intent arriving
without a valid, in-validity token.}
\reqfield{Dependency}{%
Req-T23-ACT-07}
\reqfield{Traceability (Forward)}{%
\begin{itemize}
\tightlist
\item
  Token validated at the last mile by the ACT runtime.
\end{itemize}}
\reqfield{UNITY-6G components}{%
\begin{itemize}
\tightlist
\item
  ACT runtime
\item
  Conflict fabric
\item
  DMO
\end{itemize}}
\reqfield{Linked KPIs}{%
\begin{itemize}
\tightlist
\item
  KPI\#21: Reduce the performance penalty due to conflicts by 20\% by
  means of DLT, hierarchical algorithms and policy-based conflict
  resolution schemes.
\item
  KPI\#36: Guarantee 99.99999\% service availability and reliability
  based on effective monitoring and LCM.
\end{itemize}}
\end{requirement}

\begin{requirement}{Req-T23-CG-08}{MUST}{Agents never escalate}
\reqfield{Description / rationale}{%
Type 1 conflicts shall resolve at the DMO, Type 2 shall escalate from
DMO to IDMO, and Type 3 shall be negotiated between IDMO and DMO. An
agent shall not initiate escalation and shall not address the IDMO; it
shall submit locally and receive one verdict. Per-domain members of a
federated agent shall be mediated independently in their own domains.}
\reqfield{Dependency}{%
Req-T23-DE-10}
\reqfield{Traceability (Forward)}{%
\begin{itemize}
\tightlist
\item
  Escalation paths held by the DMO and IDMO, not exposed to agents.
\end{itemize}}
\reqfield{UNITY-6G components}{%
\begin{itemize}
\tightlist
\item
  DMO
\item
  IDMO
\item
  SMB
\end{itemize}}
\reqfield{Linked KPIs}{%
\begin{itemize}
\tightlist
\item
  KPI\#26: Increase the ratio of service management and LCM tasks
  resolved by local AE/DE components.
\item
  KPI\#21: Reduce the performance penalty due to conflicts by 20\% by
  means of DLT, hierarchical algorithms and policy-based conflict
  resolution schemes.
\end{itemize}}
\end{requirement}

\begin{requirement}{Req-T23-CG-09}{MUST}{The registry is the source of authority truth}
\reqfield{Description / rationale}{%
The fabric shall resolve against the registry and never against claims
carried in a submitted message. A submission whose control parameters
exceed the registered control scope shall be rejected as a conformance
fault and shall not be arbitrated.}
\reqfield{Dependency}{%
Req-T23-LC-05}
\reqfield{Traceability (Forward)}{%
\begin{itemize}
\tightlist
\item
  Registry consulted per submission by the fabric.
\end{itemize}}
\reqfield{UNITY-6G components}{%
\begin{itemize}
\tightlist
\item
  DMO
\item
  Conflict fabric
\end{itemize}}
\reqfield{Linked KPIs}{%
\begin{itemize}
\tightlist
\item
  KPI\#36: Guarantee 99.99999\% service availability and reliability
  based on effective monitoring and LCM.
\item
  KPI\#21: Reduce the performance penalty due to conflicts by 20\% by
  means of DLT, hierarchical algorithms and policy-based conflict
  resolution schemes.
\end{itemize}}
\end{requirement}

\begin{requirement}{Req-T23-CG-10}{MUST}{Conflict outcomes are platform-emitted evidence}
\reqfield{Description / rationale}{%
The platform shall collect, per agent and per service, the verdict
distribution, unresolved conflict count, escalation count,
implicit-conflict attributions and oscillation indicators, and shall
route them to assurance, to evaluation-led authority downgrade and to
KPI\#21 measurement. These counters shall be platform-emitted and never
agent-reported.}
\reqfield{Dependency}{%
Req-T23-LC-08}
\reqfield{Traceability (Forward)}{%
\begin{itemize}
\tightlist
\item
  Counters emitted by the fabric and the platform runtime.
\end{itemize}}
\reqfield{UNITY-6G components}{%
\begin{itemize}
\tightlist
\item
  DMO
\item
  Evaluation
\item
  WP4 resolver
\end{itemize}}
\reqfield{Linked KPIs}{%
\begin{itemize}
\tightlist
\item
  KPI\#21: Reduce the performance penalty due to conflicts by 20\% by
  means of DLT, hierarchical algorithms and policy-based conflict
  resolution schemes.
\item
  KPI\#36: Guarantee 99.99999\% service availability and reliability
  based on effective monitoring and LCM.
\end{itemize}}
\end{requirement}

\begin{requirement}{Req-T23-CG-11}{MUST}{Priority claims are bounded by the registry}
\reqfield{Description / rationale}{%
An agent may express urgency, but its effective arbitration priority
shall be the one registered at admission from the service tier.
Self-declared urgency shall order intents only within the
agent\textquotesingle s own registered priority band.}
\reqfield{Dependency}{%
Req-T23-CG-09}
\reqfield{Traceability (Forward)}{%
\begin{itemize}
\tightlist
\item
  Effective priority resolved from the registry at arbitration.
\end{itemize}}
\reqfield{UNITY-6G components}{%
\begin{itemize}
\tightlist
\item
  DMO
\item
  Conflict fabric
\item
  WP4 resolver
\end{itemize}}
\reqfield{Linked KPIs}{%
\begin{itemize}
\tightlist
\item
  KPI\#33: Support more tenants via distributed AI \& DLT.
\item
  KPI\#21: Reduce the performance penalty due to conflicts by 20\% by
  means of DLT, hierarchical algorithms and policy-based conflict
  resolution schemes.
\end{itemize}}
\end{requirement}

\begin{requirement}{Req-T23-CG-12}{MUST}{Twin validation is resolver-side}
\reqfield{Description / rationale}{%
Where a network digital twin is used to simulate candidate actions, that
simulation shall be a capability of the resolver. An agent shall not
assert that its action was twin-validated, and no authority shall be
granted on the strength of such a claim. Twin validation shall
strengthen a verdict and shall never alter an agent\textquotesingle s
authority state.}
\reqfield{Dependency}{%
Req-T23-CG-09}
\reqfield{Traceability (Forward)}{%
\begin{itemize}
\tightlist
\item
  Twin invoked resolver-side; agent claims carry no weight.
\end{itemize}}
\reqfield{UNITY-6G components}{%
\begin{itemize}
\tightlist
\item
  WP4 resolver
\item
  Conflict fabric
\end{itemize}}
\reqfield{Linked KPIs}{%
\begin{itemize}
\tightlist
\item
  KPI\#21: Reduce the performance penalty due to conflicts by 20\% by
  means of DLT, hierarchical algorithms and policy-based conflict
  resolution schemes.
\item
  KPI\#42: DT-validated AI algorithms that do not break the
  time-sensitive operation when deployed in 99\% of the outputs.
\end{itemize}}
\end{requirement}

\hypertarget{contribution-to-the-project-kpis-6}{%
\subsection{CONTRIBUTION TO THE PROJECT KPIs}\label{contribution-to-the-project-kpis-6}}

Table 10-5 states what conflict resolution and governance contribute to
each KPI they bear on, and Table 10-6 maps the requirements of this
chapter onto those KPIs.

\begin{wptable}{10-5}

\begin{longtable}[]{@{}
  >{\raggedright\arraybackslash}p{(\columnwidth - 4\tabcolsep) * \real{0.1499}}
  >{\raggedright\arraybackslash}p{(\columnwidth - 4\tabcolsep) * \real{0.6375}}
  >{\raggedright\arraybackslash}p{(\columnwidth - 4\tabcolsep) * \real{0.2126}}@{}}
\caption{Contribution of conflict resolution and governance to the
project KPIs}\tabularnewline
\toprule\noalign{}
\begin{minipage}[b]{\linewidth}\raggedright
\textbf{KPI}
\end{minipage} & \begin{minipage}[b]{\linewidth}\raggedright
\textbf{Contribution}
\end{minipage} & \begin{minipage}[b]{\linewidth}\raggedright
\textbf{Requirements}
\end{minipage} \\
\midrule\noalign{}
\endfirsthead
\toprule\noalign{}
\begin{minipage}[b]{\linewidth}\raggedright
\textbf{KPI}
\end{minipage} & \begin{minipage}[b]{\linewidth}\raggedright
\textbf{Contribution}
\end{minipage} & \begin{minipage}[b]{\linewidth}\raggedright
\textbf{Requirements}
\end{minipage} \\
\midrule\noalign{}
\endhead
\bottomrule\noalign{}
\endlastfoot
\textbf{KPI\#21} & A fixed four-verdict surface lets any WP4 resolution
algorithm be substituted without touching an admitted agent, and the
same counters quantify both agent behavior and fabric effectiveness. &
Req-T23-CG-02, Req-T23-CG-03, Req-T23-CG-10 \\
\textbf{KPI\#4} & Screening and stored bounded approval keep mediation
off the ordinary cycle, so the fabric is compatible with the tightest
latency classes. & Req-T23-CG-05, Req-T23-CG-06 \\
\textbf{KPI\#23} & Screening and stored bounded approval mean a
mediation exchange takes place only where a candidate needs a verdict,
rather than once per decision. & Req-T23-CG-05, Req-T23-CG-06 \\
\textbf{KPI\#26} & Per-domain members of a federated agent are mediated
independently, and cross-domain incoherence surfaces as a Type 2
conflict rather than requiring cross-subscription. & Req-T23-CG-08 \\
\textbf{KPI\#36} & Out-of-scope submissions are conformance faults
rather than arbitration inputs, so an agent cannot enlarge its authority
by asking, and the conflict-resolution token makes granted authority an
artifact checked at execution rather than an assumption carried down the
pipeline. & Req-T23-CG-07, Req-T23-CG-09, Req-T23-CG-10 \\
\textbf{KPI\#33} & Registry-bounded priority keeps a priority policy
meaningful when many tenants contribute agents, which is what lets the
platform admit more of them without priority-based resolution collapsing
into first-come-first-served. & Req-T23-CG-11 \\
\end{longtable}

\end{wptable}

\begin{wptable}{10-6}\footnotesize

\begin{longtable}[]{@{}
  >{\raggedright\arraybackslash}p{(\columnwidth - 16\tabcolsep) * \real{0.1813}}
  >{\raggedright\arraybackslash}p{(\columnwidth - 16\tabcolsep) * \real{0.1023}}
  >{\raggedright\arraybackslash}p{(\columnwidth - 16\tabcolsep) * \real{0.1023}}
  >{\raggedright\arraybackslash}p{(\columnwidth - 16\tabcolsep) * \real{0.1023}}
  >{\raggedright\arraybackslash}p{(\columnwidth - 16\tabcolsep) * \real{0.1023}}
  >{\raggedright\arraybackslash}p{(\columnwidth - 16\tabcolsep) * \real{0.1023}}
  >{\raggedright\arraybackslash}p{(\columnwidth - 16\tabcolsep) * \real{0.1023}}
  >{\raggedright\arraybackslash}p{(\columnwidth - 16\tabcolsep) * \real{0.1023}}
  >{\raggedright\arraybackslash}p{(\columnwidth - 16\tabcolsep) * \real{0.1025}}@{}}
\caption{Mapping of conflict resolution and governance requirements onto
the project KPIs}\tabularnewline
\toprule\noalign{}
\begin{minipage}[b]{\linewidth}\raggedright
\textbf{KPI\textbackslash{}}

\textbf{Requirement}
\end{minipage} & \begin{minipage}[b]{\linewidth}\raggedright
\textbf{\#1}
\end{minipage} & \begin{minipage}[b]{\linewidth}\raggedright
\textbf{\#4}
\end{minipage} & \begin{minipage}[b]{\linewidth}\raggedright
\textbf{\#21}
\end{minipage} & \begin{minipage}[b]{\linewidth}\raggedright
\textbf{\#23}
\end{minipage} & \begin{minipage}[b]{\linewidth}\raggedright
\textbf{\#26}
\end{minipage} & \begin{minipage}[b]{\linewidth}\raggedright
\textbf{\#33}
\end{minipage} & \begin{minipage}[b]{\linewidth}\raggedright
\textbf{\#36}
\end{minipage} & \begin{minipage}[b]{\linewidth}\raggedright
\textbf{\#42}
\end{minipage} \\
\midrule\noalign{}
\endfirsthead
\toprule\noalign{}
\begin{minipage}[b]{\linewidth}\raggedright
\textbf{KPI\textbackslash{}}

\textbf{Requirement}
\end{minipage} & \begin{minipage}[b]{\linewidth}\raggedright
\textbf{\#1}
\end{minipage} & \begin{minipage}[b]{\linewidth}\raggedright
\textbf{\#4}
\end{minipage} & \begin{minipage}[b]{\linewidth}\raggedright
\textbf{\#21}
\end{minipage} & \begin{minipage}[b]{\linewidth}\raggedright
\textbf{\#23}
\end{minipage} & \begin{minipage}[b]{\linewidth}\raggedright
\textbf{\#26}
\end{minipage} & \begin{minipage}[b]{\linewidth}\raggedright
\textbf{\#33}
\end{minipage} & \begin{minipage}[b]{\linewidth}\raggedright
\textbf{\#36}
\end{minipage} & \begin{minipage}[b]{\linewidth}\raggedright
\textbf{\#42}
\end{minipage} \\
\midrule\noalign{}
\endhead
\bottomrule\noalign{}
\endlastfoot
\textbf{Req-T23-CG-01} & & & \kpimark{} & & & & \kpimark{} & \\
\textbf{Req-T23-CG-02} & \kpimark{} & & \kpimark{} & & & & & \\
\textbf{Req-T23-CG-03} & & & \kpimark{} & & & & \kpimark{} & \\
\textbf{Req-T23-CG-04} & & & \kpimark{} & & & & \kpimark{} & \\
\textbf{Req-T23-CG-05} & & \kpimark{} & \kpimark{} & \kpimark{} & & &
& \\
\textbf{Req-T23-CG-13} & & & \kpimark{} & & & & \kpimark{} & \\
\textbf{Req-T23-CG-06} & & \kpimark{} & & \kpimark{} & & & & \\
\textbf{Req-T23-CG-07} & & & \kpimark{} & & & & \kpimark{} & \\
\textbf{Req-T23-CG-08} & & & \kpimark{} & & \kpimark{} & & & \\
\textbf{Req-T23-CG-09} & & & \kpimark{} & & & & \kpimark{} & \\
\textbf{Req-T23-CG-10} & & & \kpimark{} & & & & \kpimark{} & \\
\textbf{Req-T23-CG-11} & & & \kpimark{} & & & \kpimark{} & & \\
\textbf{Req-T23-CG-12} & & & \kpimark{} & & & & & \kpimark{} \\
\end{longtable}

\end{wptable}
